My first conversion attempt used the robot model, which combined BPA routing on the Gait2392 skeleton with unchanged human muscle parameters. MyoConverter's muscle-equilibrium checks rejected infeasible combinations of routed length and muscle parameters. I therefore separated conversion of the human reference from reconstruction of the robot routing. This distinction matters because an error caused by an artificial-actuator route should not be interpreted as a property of the original human model.

The human-model conversion also required specific repairs. Three small hip rotators per side, the quadratus femoris, gemelli, and piriformis, were disabled at model load. The converted model retains 92 actuator entries, but these six contribute no force. The model retains the OpenSim and International Society of Biomechanics coordinate convention; frame transformations and individual-actuator tests check that joint motion and torque directions have the intended signs.

The source represents patella motion through muscle insertion locations that vary with knee angle. In the converted model, limits on the joints representing these moving path points conflicted with their coupling constraints. Removing those conflicting limits allowed the coupled path points to follow their prescribed trajectories. The rectus femoris path also required a patella-tracking via point to restore its knee-extension action. Explicit collision pairs introduced unwanted contacts, so contact was restricted to the foot--ground pairs. The anatomical joint limits also had to be re-enabled; these driver joints were treated separately from the coupled path-point joints. Passive joint torques supplemented the joint behavior. These repairs explain why conversion needed checks beyond successful loading: a model could load while producing an incorrect torque direction or unintended contact force.

Comparisons with the OpenSim attachment paths and single-actuator tests checked the transferred geometry and force directions. They do not establish that every muscle reproduces the source force throughout its range or that the controller produces walking. The neural and mechanical models exchange signals at a \qty{2}{\milli\second} solver time step. Appendix~\ref{sec:app_conversion_repairs} records the individual repairs and their scope.

